%
% Capítulo 6
%
\chapter{Testes e Validação} \label{cap:testes}

A validação do sistema foi realizada de forma progressiva ao longo do desenvolvimento, acompanhando a implementação das diferentes funcionalidades da aplicação e do backend.
Esta abordagem permitiu testar isoladamente componentes específicos e, posteriormente, verificar o comportamento integrado da solução. Numa fase inicial, os testes incidiram sobretudo sobre a lógica de tratamento de disponibilidades, a criação de horários, o funcionamento da API e os fluxos principais da aplicação Android.

Na versão final, a estratégia de testes evoluiu de verificações maioritariamente manuais para a implementação de uma infraestrutura robusta de testes integrados e unitários no backend, garantindo a integridade das regras de negócio, dos endpoints da API e dos fluxos das aplicações móveis.

\section{Infraestrutura e Testes Automatizados no Backend}

Para assegurar a fiabilidade da lógica de negócio sem depender da base de dados de produção (PostgreSQL), foi montado um ecossistema de testes automatizados suportado por uma base de dados \textit{H2} em memória partilhada.

Através da classe \textit{TestDatabase}, implementou-se um mecanismo que executa um \textit{reset} integral ao esquema de tabelas antes da execução de cada caso de teste. Esta abordagem garante o isolamento completo dos testes, permitindo que corram de forma rápida e previsível.

Desta forma, os testes automatizados cobrem atualmente os seguintes componentes lógicos e de serviço:

\begin{itemize}
    \item \textbf{Processador de intervalos temporais:} Responsável por dividir disponibilidades em blocos utilizáveis pelo algoritmo. Os testes unitários permitem verificar aspetos como: o número de blocos criados para diferentes intervalos, a preservação correta do dia da semana, o comportamento em intervalos inválidos ou demasiado curtos, e o suporte a diferentes durações de bloco.

    \item \textbf{Mecanismo e algoritmo de criação de horários:} Responsável por construir as propostas de sessões. Os testes unitários e de serviço (\textit{ScheduleServiceTest}) validam a atribuição correta de alunos disponíveis, o respeito pelo número máximo de participantes por sessão, o respeito pelo número máximo de sessões por aluno por dia, as restrições semanais de carga horária do aluno (\textit{weeklyHours}) e o comportamento em cenários sem alunos ou sem disponibilidade do professor.

    \item \textbf{Gestão de Aulas Concretas (\textit{LessonServiceTest}):} Uma adição crítica da versão final que testa a lógica de persistência de aulas recorrentes com datas reais, a consulta de histórico, a marcação de presenças (\textit{attendance}), os fluxos de cancelamento de séries (com simulação de notificações) e a deteção automática de conflitos em edições de horário.

    \item \textbf{Rotas e Contratos da API (\textit{Integration Level}):} Recorrendo à funcionalidade \textit{testApplication} do \textit{Ktor}, foram criados testes de integração automatizados (\textit{TeacherRoutesTest}, \textit{StudentRoutesTest}, \textit{AvailabilityRoutesTest}, \textit{RestrictionsRoutesTest}, \textit{ScheduleRoutesTest}, \textit{LessonRoutesTest}, \textit{AccountRoutesTest} e \textit{AuthRoutesExtraTest}) que simulam pedidos HTTP contra o servidor e validam os payloads JSON e os códigos de resposta face ao esperado.
\end{itemize}

Estes testes constituem uma base importante de confiança para o comportamento do sistema, sobretudo nas partes em que a lógica é mais sensível a regras e combinações de dados.

\section{Validação do Backend com Postman}

Com a introdução dos testes automatizados via \textit{Ktor} e \textit{H2}, a ferramenta Postman deixou de ser a única forma de testar a API, transitando para um papel de validação manual durante o desenvolvimento concorrente das interfaces móveis.

Esta ferramenta revelou-se importante para validar em tempo real o comportamento físico do \textit{PostgreSQL} e inspecionar os payloads de rotas mais complexas antes da sua integração final na aplicação.
Os testes realizados no Postman integraram os fluxos originais e as novas rotas implementadas, cobrindo os seguintes cenários do domínio:

\begin{itemize}
    \item verificação de disponibilidade do servidor (\textit{/health});
    \item criação de professores e alunos;
    \item autenticação de utilizadores e validação de credenciais (\textit{/login});
    \item consulta de listagens globais de professores e alunos;
    \item associação de alunos a professores e remoção desse vínculo;
    \item consulta de alunos associados a um professor;
    \item \item registo, consulta e remoção de lote de disponibilidades;
    \item registo, consulta e atualização de restrições do professor;
    \item criação e execução do algoritmo de criação de horários (\textit{POST /lessons/generate});
    \item consulta de horários semanais e históricos específicos de professores e alunos;
    \item alteração dinâmica do estado de presença de um aluno;
    \item remoção de lote de séries de aulas por \textit{seriesId} e cancelamentos individuais;
    \item envio de alertas livres por e-mail aos alunos via \textit{EmailService}.
\end{itemize}

Através destes testes foi possível confirmar que os métodos HTTP e os principais fluxos do domínio estavam corretamente representados no backend. Esta validação em duplicado permitiu também identificar e corrigir inconsistências entre o modelo de dados do frontend e o do backend, contribuindo para uma integração mais robusta entre os dois componentes do ecossistema.


\section{Testes Funcionais da Aplicação Android}

Ao nível da aplicação Android, os testes realizados incidiram sobre os principais fluxos funcionais disponíveis para cada tipo de utilizador.
O objetivo foi validar o comportamento da interface, a navegação entre ecrãs, a persistência de dados e a integração progressiva com a API.

Entre os fluxos testados destacam-se:
\begin{itemize}
    \item autenticação e gestão de sessão;
    \item distinção de comportamento entre professor e aluno;
    \item seleção e troca de professor por parte do aluno;
    \item visualização de alunos associados ao professor;
    \item definição de disponibilidades, incluindo múltiplos intervalos por dia;
    \item definição de restrições;
    \item criação e apresentação do horário;
    \item aceitação do horário e integração com Google Calendar.
    \item consulta de aulas persistidas e histórico;
    \item marcação de presenças por parte do professor;
    \item remarcação e cancelamento de aulas;
    \item envio manual de notificações aos alunos.
\end{itemize}

A validação funcional da aplicação foi importante para garantir que a interface suportava corretamente os comportamentos previstos pelo domínio
do problema e que a informação apresentada ao utilizador era consistente com o estado armazenado localmente e, quando aplicável, com a informação obtida do backend.
Foi também validada a integração com o Google Calendar no contexto do professor.
Após a aceitação de uma proposta de horário, verificou-se que as sessões podiam ser registadas com sucesso no calendário do utilizador autenticado, confirmando o funcionamento do fluxo de autenticação Google e da utilização da respetiva API.

\section{Validação da Criação de Horários}

A criação de horários foi alvo de validação específica, dado tratar-se da funcionalidade central do projeto.
Para esse efeito, foram construídos cenários de teste com diferentes combinações de disponibilidades, alunos e restrições,
permitindo observar o comportamento do sistema em situações representativas.

Os testes incidiram sobre aspetos como:
\begin{itemize}
    \item divisão correta dos intervalos de disponibilidade em blocos temporais;
    \item criação de sessões compatíveis com a disponibilidade do professor;
    \item alocação de alunos apenas a sessões em que se encontram disponíveis;
    \item respeito pelos limites de participantes por sessão;
    \item respeito pelo número máximo de sessões por aluno no mesmo dia;
    \item distribuição de alunos por diferentes blocos quando aplicável.
    \item respeito pelo limite semanal de horas definido para cada aluno;
\end{itemize}

Através destes testes foi possível confirmar que o projeto já produz propostas de horário coerentes com as regras definidas e com a informação introduzida no sistema.
Embora a abordagem atual ainda possa ser melhorada, os resultados obtidos mostram que a funcionalidade principal do projeto já se encontra operacional.

\section{Integração e Validação Progressiva}

Um aspeto relevante da fase de testes foi a validação progressiva da integração entre frontend e backend. Em vez de tentar validar toda a solução apenas no final,
as funcionalidades foram sendo testadas à medida que a comunicação remota era introduzida.
Esta estratégia permitiu identificar mais cedo inconsistências em endpoints, modelos de dados, fluxos de autenticação e mecanismos de sincronização.

A coexistência de persistência local e persistência remota exigiu também atenção adicional, uma vez que parte da validação teve de assegurar que os dados obtidos da API
se mantinham consistentes com o estado visível na aplicação.
Esta abordagem incremental revelou-se útil para consolidar a solução final e estabelecer uma base sólida para futuras evoluções funcionais do sistema.

\section{Síntese dos Resultados Obtidos}

Os testes realizados permitiram concluir que o sistema apresenta um nível de funcionalidade significativo. O backend está operacional, os principais endpoints foram validados com sucesso, a aplicação Android executa os fluxos centrais do sistema, a autenticação encontra-se funcional e a criação de horários produz resultados coerentes em cenários representativos.

Apesar de persistirem aspetos a melhorar, nomeadamente ao nível da integração total entre componentes, da evolução futura do algoritmo e do polimento de algumas componentes da interface, a validação efetuada demonstra que a solução
já consegue suportar os casos de uso mais importantes do projeto.
Este resultado constitui um indicador positivo da robustez da solução atualmente desenvolvida e da sua capacidade de suportar futuras evoluções.